Memory for LLM agents is usually built as retrieval over an append-only store, with an LLM adjudicating conflicts at query time. This works until the evidence itself changes. Figure~\ref{fig:example} shows the case this paper is about. A press report says Alex works for Veltran; a registry says Veltran's head office is in Tessaly; a rule the agent is allowed to use says a person works where their employer is headquartered, so the agent concludes that Alex works in Tessaly and, from a second rule, that Alex pays local tax there. On day 12 the press withdraws its report. Nothing about Veltran changed and no new report about Alex arrived, so a store that updates beliefs only when new values are asserted keeps saying Tessaly; a store that re-derives everything from scratch at query time gets it right but pays for the re-derivation at every question; and an LLM asked to adjudicate the whole log must notice that a withdrawal two steps upstream has removed the only justification for the answer it is about to give. The question is narrower than ``how should agents learn continually'' and, we argue, prior to it: \emph{when evidence is withdrawn or corrected, which conclusions should lose support, and what does it cost to maintain that structure at write time rather than re-derive it at query time?}

\begin{figure}[t]\centering
\begin{tikzpicture}[>=stealth, node distance=4mm and 9mm, every node/.style={font=\footnotesize}, obs/.style={draw, rounded corners=1pt, inner sep=2.5pt, fill=blue!4}, der/.style={draw, inner sep=2.5pt, fill=orange!8}, lab/.style={font=\scriptsize, text=black!60}]
\node[obs] (o1) {o1 (press, day 3): employer(Alex) = Veltran};
\node[obs, below=of o1] (o2) {o2 (registry, day 5): hq\_city(Veltran) = Tessaly};
\node[der, right=of $(o1.east)!0.5!(o2.east)$, xshift=6mm] (w) {work\_city(Alex) = Tessaly};
\node[der, right=of w] (t) {local\_tax\_city(Alex) = Tessaly};
\draw[->] (o1.east) -- (w.west); \draw[->] (o2.east) -- (w.west); \draw[->] (w) -- (t);
\node[obs, fill=red!8, below=of o2] (r) {o3 (press, day 12): \emph{withdraws} o1};
\draw[->, red!70!black, dashed] (r.east) to[out=0, in=250] (w.south);
\node[lab, below=1mm of t, align=left] {day 13: justified answer is \texttt{unknown};\\ a value-triggered store still says Tessaly};
\end{tikzpicture}
\caption{A retraction consequence. Derived beliefs are justified by the reports they were derived from (solid arrows). The withdrawal of o1 removes the only justification of both derived beliefs although no report about them changed. Scoring is against what the surviving evidence \emph{justifies}, not against the world.}\label{fig:example}
\end{figure}

Three layers have to be kept apart to study this. The \emph{world} is what is actually the case; it is hidden from every system and used only to generate reports. The \emph{observations} are what sources reported, when, with what reliability class, and which reports were later withdrawn or corrected. \emph{Justified belief} is what the observations warrant under declared assumptions: not the world (a true belief held for no good reason earns nothing), and not the latest report (a later statement is evidence, not truth). Systems are scored against the third layer. Two design commitments follow. The reference reasoner sees only the observations and the declared assumptions, never the world, and when the evidence does not decide between readings it \emph{retains} them rather than selecting one (\S\ref{sec:semantics}). And the requirement to preserve that uncertainty is placed on observable behaviour through a shared output contract (\S\ref{sec:contract}), so that a system representing hypotheses explicitly and an LLM adjudicating a flat log compete on equal terms.

What the study shows. (i)~A declared semantics of justified belief for observation streams with operator cues, dispute, correction and withdrawal, whose per-key implementation is validated against a brute-force global enumeration on \ovQueries{} value queries (Appendix~\ref{app:semantics}), and a benchmark generator, \rs, that separates the three layers with a declared menu of perturbations. (ii)~\pal, an incremental store in which derived beliefs pin the versions of the evidence that justified them, returns the same status, value and alternatives as a symbolic replay of the same semantics on every query of every setting (after one bug fix we report; provenance lists differ on 2--5\%), at six to twelve times lower measured query latency; its total work is lower than replay only when queries are at least as frequent as observations or histories are long, and its write cost grows exponentially in reports per key (\S\ref{sec:res-systems}). (iii)~Within that implementation, retraction propagation through justifications, not ordinary belief updating, is what prevents downstream errors: switching it off costs about 20 points on warranted downstream queries in every setting and nothing on direct ones. (iv)~A regime analysis against the hidden world quantifies what refusing to assume recency buys and costs: last-write-wins is right on about 92\% of value queries whenever the latest report is true (\rgLatestTruePct\% of the queries in our generator) and on 3--12\% when it is delayed or erroneous; justified belief abstains on 20--45\% of the former and on 59--94\% of the latter; a pre-registered same-origin self-update rule removes the abstention in the single-source case at a cost of at most 1.3 points of wrong commitments where the latest report is wrong (\S\ref{sec:res-regime}). (v)~Under automatically extracted assertions with a fixed, declared rule schema, symbolic justification maintenance beats a prompted LLM adjudicator over the same claims by \sTwoHOne{} points on warranted downstream queries while committing on about a seventh as many unresolved slots; a rule-based adjudicator over the same claims scores the same as \pal, so this is a result about the evaluated adjudicators, not about write time versus query time (\S\ref{sec:res-s2}). (vi)~One pre-registered hypothesis fails (calibration) and one could not be executed as registered (the matched-cost clause of H1); the corruption curve locates the error rates at which last-write-wins overtakes justified belief; and when the latest report is an attacker's injection the store commits to it on \rgPoisPalInj{} of queries against \rgPoisLwwInj{} for recency.

Retraction propagation as a \emph{mechanism} is not new: incremental view maintenance keeps derived relations consistent under base-table changes \citep{gupta1993}, justification- and assumption-based truth maintenance withdraw conclusions whose premises are withdrawn \citep{doyle1979,dekleer1986}, MemTX cascades typed repair through derived records in transactional agent memory \citep{memtx2026}, and CUPMem invalidates dependent attributes at write time \citep{stale2026}. Our contribution is the measurement: a semantics and a scoring contract under which premature commitment is visible, a benchmark in which the three layers are separable, and the operating regimes, costs and failure modes of maintaining justifications, including the ones that went against the method.
